Finiteness and termination of the symbolic construction are \(\text{thm:tcf-termination}\) applied to the corrected history-preserving procedure and budget. The construction of \(\mathscr P_{\mathcal G,A,Y}\) in \(\text{def:symbolic-observed-partition}\), together with \(\text{lem:semialgebraic-fiber-dichotomy}\), gives a finite, pairwise-disjoint partition of the feasible projection and ensures that each compatible observed law belongs to exactly one cell.

Fix such a law, \(a\in\mathcal X_A\), and \(y\in\mathcal X_Y\). The dichotomy lemma and both directions of \(\text{thm:response-fiber-equivalence}\), which preserve \(q_{a,y}(\theta)=Q_a^{\mathcal M}(y)\), give
\[
\begin{aligned}
&\text{the cell has a unique value for }q_{a,y}\\
&\quad\Longleftrightarrow
q_{a,y}\text{ is constant on }
\{\theta:\mathsf C_{\mathcal G}(\theta;P_{\mathrm{obs}})\}\\
&\quad\Longleftrightarrow
\mathcal M\longmapsto Q_a^{\mathcal M}(y)
\text{ is constant on }\mathfrak F(\mathcal G,P_{\mathrm{obs}}).
\end{aligned}
\tag{1}
\]
On the true side of (1), let \(z=f_a(P_{\mathrm{obs}})(y)\) be the unique value in \(\mathsf V_{a,y}(P_{\mathrm{obs}},z)\). For any \(\mathcal M\in\mathfrak F(\mathcal G,P_{\mathrm{obs}})\), the forward direction of \(\text{thm:response-fiber-equivalence}\) supplies \(\theta_{\mathcal M}\) satisfying
\[
\mathsf C_{\mathcal G}(\theta_{\mathcal M};P_{\mathrm{obs}}),
\qquad
q_{a,y}(\theta_{\mathcal M})=Q_a^{\mathcal M}(y).
\]
Consequently \(\mathsf V_{a,y}(P_{\mathrm{obs}},Q_a^{\mathcal M}(y))\) holds by \(\text{def:terminal-qe-test}\). Uniqueness of \(z\) therefore gives, for every model in the fiber,
\[
f_a(P_{\mathrm{obs}})(y)=Q_a^{\mathcal M}(y).
\tag{2}
\]
On the false side of (1), \(\text{lem:semialgebraic-fiber-dichotomy}\) supplies \(\mathsf W_j\). By \(\text{def:tcf-procedure}\), that relation is attached to a terminal package, and \(\text{thm:tcf-termination}\) makes its operation trace finite. Under exact specialization, \(\text{thm:qe-failure-compiler}\) applied to any extracted sample and that finite trace gives \(\mathcal M_0,\mathcal M_1\in\mathfrak F(\mathcal G,P_{\mathrm{obs}})\), strictly positive full-data cells, individual-edge-factorized latent laws, and
\[
P_{\mathcal M_0}(\mathcal O)=P_{\mathcal M_1}(\mathcal O)=P_{\mathrm{obs}},
\qquad
Q_a^{\mathcal M_0}(y)\neq Q_a^{\mathcal M_1}(y).
\tag{3}
\]

Finally suppose (2) holds for every \(y\) at \(a=a_0\) and \(a=a_1\). The finite alphabets make the following sums finite. By \(\text{def:causal-query}\) and then substitution of (2),
\[
\begin{aligned}
\tau_{s_Y}
&=\sum_{y\in\mathcal X_Y}s_Y(y)
  \{Q_{a_1}^{\mathcal M}(y)-Q_{a_0}^{\mathcal M}(y)\}\\
&=\sum_{y\in\mathcal X_Y}s_Y(y)
  \{f_{a_1}(P_{\mathrm{obs}})(y)-f_{a_0}(P_{\mathrm{obs}})(y)\}.
\end{aligned}
\tag{4}
\]
Equations (1)--(4) prove the pointwise identification and nonidentification alternatives and their relationship to the causal contrast. No step infers a colored hedge from a failed graphical branch or aligns an algebraic witness with a candidate forest, so the proof establishes exactly the narrowed statement.